\PassOptionsToPackage{dvipsnames,svgnames,table}{xcolor}
\documentclass[11pt]{article}
\usepackage[a4paper,margin=22mm]{geometry}
\usepackage[no-math]{fontspec}
\setmainfont{DejaVu Serif}
\usepackage{amsmath,amssymb,amsthm,mathtools,graphicx,xcolor,hyperref,xurl,xspace,ifthen,enumerate}
\usepackage[shortlabels]{enumitem}
\setlength{\emergencystretch}{4em}
\AtBeginDocument{\special{pdf:mapline pzdr Dingbats <uzdr.pfb}}
\SetMathAlphabet{\mathbf}{normal}{OT1}{cmr}{bx}{n}
\SetMathAlphabet{\mathbf}{bold}{OT1}{cmr}{bx}{n}
\newtheorem{theo}{Theorem}[section]
\newtheorem{lemm}[theo]{Lemma}
\newtheorem{thm}[theo]{Theorem}
\newtheorem{proposition}[theo]{Proposition}
\newtheorem{corollary}[theo]{Corollary}
\newtheorem{definition}[theo]{Definition}
\newtheorem{theorem}[theo]{Theorem}
\newtheorem*{remas*}{Remarks}
\newtheorem{remas}[theo]{Remarks}
\newtheorem*{exam*}{Example}
\newtheorem{ques}[theo]{Question}
\providecommand{\og}{«\nobreakspace}
\providecommand{\fg}{\nobreakspace»}
\newtheorem{lem}[theo]{Lemma}
\newtheorem{lemma}[theo]{Lemma}
\newtheorem{prop}[theo]{Proposition}
\newtheorem{coro}[theo]{Corollary}
\newtheorem{corol}[theo]{Corollary}
\newtheorem*{conj*}{Conjecture}
\newtheorem{conj}[theo]{Conjecture}
\newtheorem{conjecture}[theo]{Conjecture}
\newtheorem{Proposition}[theo]{Proposition}
\newtheorem{theoreme}[theo]{Theorem}
\theoremstyle{definition}
\newtheorem{defi}[theo]{Definition}
\newtheorem{exam}[theo]{Example}
\newtheorem{exem}[theo]{Example}
\newtheorem{remark}[theo]{Remark}
\newtheorem{example}[theo]{Example}
\newtheorem{exams}[theo]{Examples}
\newtheorem{rema}[theo]{Remark}
\newtheorem{nota}[theo]{Notation}
\newtheorem{notation}[theo]{Notation}
\newtheorem*{notas}{Notations}
\newtheorem{result}[theo]{Result}
\newtheorem{question}[theo]{Question}
\newtheorem{prob}[theo]{Problem}
\newtheorem*{rema*}{Remark}
\newtheorem*{defi*}{Definition}
\newtheorem*{theo*}{Theorem}
\newtheorem*{lemm*}{Lemma}
\newtheorem*{prop*}{Proposition}
\newcommand{\enoncename}{}
\newtheorem{corpusenonce}[theo]{\enoncename}
\newenvironment{enonce}[2][]{\renewcommand{\enoncename}{#2}\begin{corpusenonce}}{\end{corpusenonce}}
\newenvironment{enonce*}[2][]{\par\medskip\noindent\textbf{#2.}\itshape}{\par\medskip}
\newenvironment{proofc}[1][\proofname]{\begin{proof}[#1]}{\end{proof}}
\newcommand{\Mc}{\mathcal{M}}
\newcommand{\pointir}{.\ ---\ }
\newcommand{\equalenv}[2]{\expandafter\let\csname#1\expandafter\endcsname\csname#2\endcsname\expandafter\let\csname end#1\expandafter\endcsname\csname end#2\endcsname}
\providecommand{\diagbox}[3][]{\begin{tikzpicture}[x=1em,y=1em,baseline=(current bounding box.center)]\useasboundingbox(0,0) rectangle(3,1.4);\draw(0,1.4)--(3,0);\node at(.5,.3){#2};\node at(2.5,1.1){#3};\end{tikzpicture}}
\providecommand{\eg}{e.g.\ }
\providecommand{\ie}{i.e.\ }
\providecommand{\cf}{cf.\ }
\providecommand{\longthanks}[1]{\par\smallskip\noindent #1\par\smallskip}
\providecommand{\qbinom}[2]{\genfrac{[}{]}{0pt}{}{#1}{#2}}

\providecommand{\mathof}[1]{\mathrm{#1}}

\numberwithin{equation}{section}


\def\endappendix{}


\usepackage{latexsym,amsmath,amssymb,verbatim,tikz,tikz-cd}
\usepackage{mathtools}
\usepackage{graphicx,xcolor}

\usepackage[english]{babel}
\usepackage{amssymb}
\usepackage{amsmath}
\usepackage{amsthm}
\usepackage{graphicx}
\usepackage{csquotes}
\usepackage{bm}
\usepackage{enumitem}
\usepackage{subcaption}
\usepackage{caption}
\renewcommand{\thesubfigure}{\Alph{subfigure}}
\usepackage{dirtytalk}
\providecommand{\nopunct}{\spacefactor1007\relax}
\newtheorem{problem}[theo]{Problem}

\newcommand{\sym}{\mathrm{Sym}}
\newcommand{\qsym}{\mathrm{QSym}}
\newcommand{\Des}{\operatorname{Des}}
\newcommand{\Asc}{\operatorname{Asc}}
\newcommand{\Syt}{\operatorname{SYT}}
\newcommand{\Short}{\operatorname{Short}}
\newcommand{\Loop}{\operatorname{Loop}}
\newcommand{\Stable}{\operatorname{Stable}}
\newcommand{\len}{\operatorname{len}}
\newcommand{\height}{\operatorname{height}}
\newcommand{\core}{\operatorname{core}}
\newcommand{\res}{\operatorname{res}}
\newcommand{\row}{\operatorname{row}}
\newcommand{\insrt}{\operatorname{insert}}
\newcommand{\UP}{\operatorname{UP}}
\newcommand{\DOWN}{\operatorname{DOWN}}


\newcommand{\ZZ}{\mathbb{Z}}
\newcommand{\NN}{\mathbb{N}}
\newcommand{\RR}{\mathbb{R}}
\newcommand{\QQ}{\mathbb{Q}}
\newcommand{\PP}{\mathbb{P}}

\newcommand{\F}{\mathcal{F}}
\newcommand{\Q}{\mathcal{Q}}
\newcommand{\A}{\mathcal{A}}
\newcommand{\M}{\mathcal{M}}
\newcommand{\I}{\mathcal{I}}
\newcommand{\PM}{\mathcal{PM}}
\newcommand{\Pcal}{\mathcal{P}}


\newcommand{\odds}{\mathrm{odds}}

\makeatletter
\newcommand{\proofstep}[1]{
  \par
  \addvspace{\medskipamount}
  \textbf{#1\@addpunct{.}}\enspace\ignorespaces
}

\newcommand\restr[2]{{
  \left.\kern-\nulldelimiterspace 
  #1 
  \vphantom{\big|} 
  \right|_{#2} 
  }}
\makeatother












\usepackage{cleveref}
\crefname{originaltheorem}{original theorem}{original theorems}
\Crefname{originaltheorem}{Original theorem}{Original theorems}
\makeatletter
\newlabel{thm:criterion schur positive young}{{2.5}{source}{}{sourcepointer-1}{}}
\newlabel{thm:criterion schur positive young@cref}{{[originaltheorem][0][]2.5}{[1][0][]source}{}{}{}}
\makeatother
\title{A set of matchings closed under Knuth equivalence is Schur-positive with two-row expansion}
\author{Avichai Marmor}
\date{doi:10.5802/alco.351; CC BY 4.0}
\begin{document}
\maketitle
\noindent\textit{Editorial note.} The counted statement is the author's labelled corollary `cor:closed under knuth like is schur positive` (source lines 1214-1221); the proof environment begins on the next line after \textbackslash{}end\{corollary\} apart from one blank line, so the pairing is by adjacency.\par\medskip
\noindent\textit{Editorial note.} Two context spans are retained: the subsection label invoked by the proof and the theorem characterising Knuth equivalence, whose statement is retained so that the invoked label is defined inside the package.\par\medskip
\noindent\textit{Editorial note.} The retained preamble is the editorially prepared standalone preamble retained for this source.\par\medskip
\noindent\textit{Editorial note.} One external reference records the Schur-positivity criterion invoked at the last line; it is a navigation pointer only.\par\medskip
\noindent\textit{Editorial scope.} One original result of the pinned publisher article, counted once: Schur-positivity of sets closed under Knuth equivalence (author label cor:closed under knuth like is schur positive).\par\medskip
\noindent\textit{Allowed standard prerequisites (editorial).}\begin{itemize}
\item Knuth equivalence for matchings. Two matchings are Knuth equivalent if they can be connected by the insertion and deletion moves of the section.. The hypothesis is that the set is closed under this equivalence.
\item Short statistic and unstable chords. Short is the statistic recording the positions of the short chords, and k denotes the number of unstable chords of a matching.. The expansion is written in terms of s\_\{N-k,k\} with k read off from any matching in the class.
\item The transformation F. F(m) = (core(m), T(m)) is the bijection from matchings to pairs consisting of a core and a standard Young tableau, discussed in the retained subsection.. Its restriction to a Knuth class is the statistic-preserving bijection at the heart of the proof.
\end{itemize}\medskip
\ifdefined\hypertarget\hypertarget{sourcepointer-1}{}\fi
\noindent\textit{Source reference.} Original theorem 2.5: sources/alco-351-matching-schur-2024/publisher-source.tex, source lines 368--368. Original criterion of Adin and Roichman invoked at source line 1228. It is a navigation pointer: the counted corollary's own proof is retained line for line and nothing is replaced by this reference. The reference number was inferred by counting the theorem-class environments of that section of the pinned source; it is not confirmed by the retained bibliography or by a label-number dump of the original, so it must be read as a navigation pointer only.\par\medskip
\label{subsec:bijective proof matchings schur positive}


\begin{theorem}
\label{thm:knuth like equivalent iff same core}
    Two matchings $m_1, m_2 \in \M_N$ are Knuth equivalent if and only if $\core(m_1) = \core(m_2)$.
\end{theorem}

\begin{corollary}
\label{cor:closed under knuth like is schur positive}
    If a set $\M \subseteq \M_N$ is closed under Knuth equivalence then it is Schur-positive with respect to $\Short$. Moreover, if $\M$ is a Knuth equivalence class then its generating function is
    \[
        \Q(\M) = s_{N-k,k},
    \]
    where $k$ is the number of unstable chords of some arbitrary matching $m \in \M$.
\end{corollary}

\begin{proof}
    Clearly, a disjoint union of Schur-positive sets is Schur-positive. Therefore, it suffices to establish the statement for equivalence classes. Let $\M \subseteq \M_N$ be an equivalence class. By Theorem~\ref{thm:knuth like equivalent iff same core} there exist $0 \le k \le \frac{N}{2}$, $f \ge 0$ and $m_0 \in \M_{N-2k,f}(\Short=\emptyset)$, such that $\M = \{m \in \M_{N,f} \mid \core(m) = m_0\}$. Recall the transformation $F:m \mapsto (\core(m),T(m))$ discussed in Section~\ref{subsec:bijective proof matchings schur positive}, and consider its restriction to $\M$. We obtain that the restricted transformation
    \[
        \restr{F}{\M}:\M \to \{m_0\} \times \Syt(N-k,k)
    \]
    is a statistic-preserving bijection. Applying Theorem~\ref{thm:criterion schur positive young} completes the proof.
\end{proof}

\begin{thebibliography}{MMMMMM}
\bibitem[3]{adin_roichman_fine_set} Adin, Ron M.; Roichman, Yuval Matrices, characters and descents, Linear Algebra Appl., Volume 469 (2015), pp. 381-418 \href{https://alco.centre-mersenne.org/articles/10.5802/alco.351/#adin_roichman_fine_set}{Publisher reference}.
\bibitem[13]{knuth_classes_schur_positive} Gessel, Ira M. Multipartite P-partitions and inner products of skew Schur functions, Combinatorics and algebra (Boulder, Colo., 1983) (Contemp. Math.), Volume 34, Amer. Math. Soc., Providence, RI, 1984, pp. 289-317 \href{https://alco.centre-mersenne.org/articles/10.5802/alco.351/#knuth_classes_schur_positive}{Publisher reference}.
\bibitem[14]{gessel1993counting} Gessel, Ira M.; Reutenauer, Christophe Counting permutations with given cycle structure and descent set, J. Combin. Theory Ser. A, Volume 64 (1993) no. 2, pp. 189-215 \href{https://alco.centre-mersenne.org/articles/10.5802/alco.351/#gessel1993counting}{Publisher reference}.
\bibitem[16]{S3_patterns_list} Hamaker, Zachary; Pawlowski, Brendan; Sagan, Bruce E. Pattern avoidance and quasisymmetric functions, Algebr. Comb., Volume 3 (2020) no. 2, pp. 365-388 \href{https://alco.centre-mersenne.org/articles/10.5802/alco.351/#S3_patterns_list}{Publisher reference}.
\bibitem[26]{sagan_book} Sagan, Bruce E. The symmetric group: Representations, combinatorial algorithms, and symmetric functions, Graduate Texts in Mathematics, 203, Springer-Verlag, New York, 2001, xvi+238 pages \href{https://alco.centre-mersenne.org/articles/10.5802/alco.351/#sagan_book}{Publisher reference}.
\end{thebibliography}
\end{document}
